\section{Background and Related Work}
\label{sec:background}

\subsection{Reading Service Paths at the Network Layer}

Users reach services through domain names, anycast targets, and replica
deployments. Routing selects the resulting path, and traceroute exposes the
hops and AS transitions along it. Path-diversity measurements count these
transitions and compare them across vantage points.

Prior work measures the user's view of submarine criticality
~\cite{liu2020outofsight}, characterizes anycast deployment
~\cite{cicalese2015anycast,devries2017anycast}, and relates cable
infrastructure to AS-level structure~\cite{ye2023impact}. These measurements
describe network-layer reachability. Our question is whether diversity at that
layer corresponds to breadth across feasible physical directions.

\subsection{Mapping Paths onto Submarine Infrastructure}

Internet Atlas and InterTubes associate long-haul facilities with observed
paths~\cite{durairajan2013atlas,durairajan2015intertubes}. iGDB integrates
facilities, fiber, and network entities into a cross-layer map
~\cite{anderson2022igdb}. Nautilus associates IP links with candidate
submarine cables and attaches a confidence score to each assignment
~\cite{ramanathan2023nautilus}. Calypso combines cable layouts, network
relationships, and inland connectivity to map routes and defines Route Stress
over those assignments~\cite{wang2025calypso}. Xaminer applies cross-layer
maps to infrastructure analysis under disaster models
~\cite{ramanathan2024xaminer}. Carisimo et al. use intercontinental-link
measurements to identify gateway routers~\cite{carisimo2023hopaway}.

These systems produce an assignment from a link or path to physical
infrastructure that could plausibly carry it. Their validation evaluates those
assignments using cable failures, targeted traceroutes, operator maps, or
expert feedback. CLASP adopts this projection approach; it does not claim to
identify ground-truth cable use.

\subsection{From Mapping to Distribution Measurement}

An assignment does not by itself describe how observation mass is distributed
across physical directions or whether that distribution tracks the one visible
at the network layer. CLASP aggregates assignments into two population-level
quantities. \emph{Submarine exposure} is the fraction of valid traceroutes
containing a feasible inter-region corridor. \emph{Corridor concentration}
describes how candidate-bearing segment mass is distributed across corridors.
A country can enter the candidate space frequently and distribute observations
widely, or enter infrequently with the candidate-bearing observations
concentrated on one direction. Section~\ref{sec:framework} defines both
quantities and preserves their separate denominators.
